\Needspace{8\baselineskip}
\SourceLeaf{94}%
\noindent Die Definitionsgleichung für $\lambda'$ liefert nach kurzer Umrechnung\hfill\\
zunächst \mbox{$\bigl(\frac{\lambda'}{m'}\bigr)^{2}=c^{2}$}; hier hat man:

\[
\frac{\lambda'}{m'}=-c,\quad\text{jedoch weiter}\quad
\frac{\lambda''}{m''}=c',\qquad \frac{\lambda'''}{m'''}=c'',\ \ldots
\]

\noindent zu setzen, um die Formeln der vorliegenden Notiz zu gewinnen. Folglich bestehen die
Gleichungen.

\[
\nu=-u,\qquad P=\frac{(1+t)^{2}}{4t}\,\frac{1-u}{a^{(\infty)}},\qquad
Q=\frac{(1-t)^{2}}{4t}\,\frac{1+u}{a^{(\infty)}},
\]

\noindent wie aus den Schlussformeln in Art.\,19 a.\,a.\,O.\ hervorgeht. Durch Addition und
Subtraction dieser letzten beiden Formeln findet man:

\[
A^{(0)}=\frac{h-u}{a^{(\infty)}},\qquad A^{(1)}=\frac{1-hu}{a^{(\infty)}}.
\]

Für’s zweite subsumirt sich die hier gelöste Aufgabe der Entwicklung von
$(h-\cos\varphi)^{-\frac{3}{2}}$ dem allgemeinen in der Dissertatiuncula (p.\,35 ff.)
behandelten Probleme. Es werden daselbst unter Nr.\,7 die dreigliedrigen \textsc{Euler}schen
Relationen aufgestellt, vermöge deren man aus den beiden ersten Coefficienten
$A^{(0)}$, $A^{(1)}$ successive alle folgenden $A^{(2)}$,
$A^{(3)}$, $\ldots$\ berechnen kann. Für die Quotienten $B^{(k)}$ der
$A$ entspringen aus den \textsc{Euler}schen Recursionsformeln die hier von \textsc{Gauss} aufgestellten
Formeln zwischen $B^{(k)}$ und $B^{(k+1)}$.

\Needspace{20\baselineskip}
Es ist endlich noch erwähnenswerth, dass sich der für $B^{(k)}$ angegebene
Kettenbruch direct der allgemeinen Formel [25] in Art.\,12 der »Disquisitiones
generales circa seriem infinitam etc.« (ges. Werke III p.\,134) unterordnet, in welcher
der Quotient zweier hypergeometrischen Functionen in einen Kettenbruch entwickelt wird.
In den daselbst erklärten Bezeichnungen würde sich die Formel für $B^{(k)}$
folgendermassen schreiben:

\[
B^{(k)}=\frac{k+\frac{3}{2}}{k+1}\,t\mathord{.} G\bigl(-\tfrac{1}{2},k-\tfrac{1}{2},k+1;t^{2}\bigr)
\]

\noindent oder auch:

\[
B^{(k)}=\frac{k+\frac{3}{2}}{k+1}\,t\mathord{.}
\frac{F\bigl(-\tfrac{1}{2},k+\tfrac{1}{2},k+2;t^{2}\bigr)}{F\bigl(-\tfrac{1}{2},k-\tfrac{1}{2},k+1;t^{2}\bigr)}.
\]

\begin{flushright}
\textsc{Fricke.}
\end{flushright}
